\documentclass[11pt,openany]{book}

\usepackage{fontspec}
\usepackage{geometry}
\geometry{paper=a4paper, margin=1in}
\usepackage{xcolor}
\usepackage{graphicx}
\usepackage{booktabs}
\usepackage{tabularx}
\usepackage{array}
\usepackage{longtable}
\usepackage{enumitem}
\usepackage{hyperref}
\usepackage{xurl}
\usepackage{titlesec}
\usepackage{fancyhdr}
\usepackage{microtype}
\usepackage{setspace}
\usepackage[most]{tcolorbox}

\setmainfont{TeX Gyre Pagella}
\setsansfont{TeX Gyre Heros}
\setmonofont{DejaVu Sans Mono}
\setstretch{1.08}

\definecolor{ink}{HTML}{14213D}
\definecolor{gold}{HTML}{B98A2F}
\definecolor{paper}{HTML}{FBF8F1}
\definecolor{mist}{HTML}{E8EEF5}
\definecolor{signal}{HTML}{F4E8C8}
\definecolor{line}{HTML}{D7D1C3}

\hypersetup{
  colorlinks=true,
  linkcolor=ink,
  urlcolor=ink,
  pdftitle={Wealth From First Principles},
  pdfauthor={LazyingArt LLC}
}

\titleformat{\chapter}[display]
  {\normalfont\sffamily\bfseries\Huge\color{ink}}
  {\filleft\Large\color{gold}\thechapter}
  {1ex}
  {\titlerule[1pt]\vspace{1ex}\filright}

\titleformat{\section}
  {\normalfont\sffamily\bfseries\Large\color{ink}}
  {\thesection}
  {0.75em}
  {}

\pagestyle{fancy}
\fancyhf{}
\fancyhead[LE,RO]{\sffamily\small\thepage}
\fancyhead[RE]{\sffamily\small Wealth From First Principles}
\fancyhead[LO]{\sffamily\small LazyingArt LLC}
\setlength{\headheight}{14pt}

\setlist[itemize]{leftmargin=1.4em}
\setlist[enumerate]{leftmargin=1.6em}
\setlength{\parindent}{0pt}
\setlength{\parskip}{0.75em}

\newtcolorbox{signalbox}{
  colback=signal,
  colframe=gold,
  boxrule=0.6pt,
  arc=2pt,
  left=10pt,
  right=10pt,
  top=8pt,
  bottom=8pt
}

\newtcolorbox{notebox}{
  colback=mist,
  colframe=ink,
  boxrule=0.6pt,
  arc=2pt,
  left=10pt,
  right=10pt,
  top=8pt,
  bottom=8pt
}

\begin{document}
\frontmatter

\begin{titlepage}
\pagecolor{paper}
\color{ink}
\vspace*{2.8cm}
{\sffamily\Huge\bfseries Wealth From First Principles\par}
\vspace{0.4cm}
{\Large A practical field guide to money, value, ownership, and financial independence\par}
\vspace{1.4cm}
\begin{signalbox}
\textbf{Core thesis}\par
Money is a claim on resources. Wealth is durable control over assets, cash flow, capabilities, and time. Most lasting wealth comes from ownership, compounding, and staying in the game long enough for good decisions to matter.
\end{signalbox}
\vfill
{\large LazyingArt LLC\par}
{\large April 2026\par}
\end{titlepage}
\nopagecolor

\tableofcontents

\mainmatter

\chapter{The shortest useful answers}

This book is organized around a cluster of recurring questions:

\begin{itemize}
  \item What is money?
  \item Where does it come from?
  \item What is wealth?
  \item Who can build it?
  \item Why do some people build more of it than others?
  \item What methods actually work in practice?
\end{itemize}

\begin{center}
\begin{tabularx}{\textwidth}{>{\raggedright\arraybackslash}p{0.29\textwidth}X}
\toprule
\textbf{Question} & \textbf{Working answer} \\
\midrule
What is money? & A shared accounting and settlement tool used to price things, exchange claims, and move purchasing power through time. \\
Where does money come from? & At the system level, mostly from bank lending, central bank liabilities, government spending, and external flows. At the personal level, from labor, business revenue, investment income, transfers, and asset sales. \\
What is wealth? & Assets minus liabilities, plus the practical ability of those assets and capabilities to keep producing value later. \\
Who can build wealth? & Many people can improve their position, but not from identical starting points and not with identical strategies. \\
Why do some people get richer faster? & They combine ownership, skill, leverage, compounding time, institutional support, and luck. \\
\bottomrule
\end{tabularx}
\end{center}

\begin{signalbox}
\textbf{A sentence worth remembering}\par
Money is a flow tool. Wealth is a stock of productive claims.
\end{signalbox}

\chapter{What money is}

Money is easy to use and surprisingly hard to define cleanly. The practical definition is better than the metaphysical one: money is whatever a society widely accepts for pricing, payment, settlement, and the storage of near-term purchasing power.

\section{The three classic functions}

\begin{enumerate}
  \item \textbf{Unit of account}. Wages, rents, profits, taxes, and debts are stated in money.
  \item \textbf{Medium of exchange}. Money removes the need for barter.
  \item \textbf{Store of value}. Money lets people move purchasing power over time, though inflation weakens this function.
\end{enumerate}

In ordinary life, the money people use most is not paper currency. It is bank deposits. This is why public debates about money often go wrong: people imagine cash, while modern economies mainly transact with digitally recorded bank liabilities.

\section{What money is not}

Money is not the same thing as income, profit, capital, or wealth.

\begin{itemize}
  \item \textbf{Income} is what arrives over a period.
  \item \textbf{Profit} is what remains after costs.
  \item \textbf{Capital} is a productive asset base.
  \item \textbf{Wealth} is accumulated net value and durable capability.
\end{itemize}

Someone can have large money inflows and little wealth if they retain no ownership and save none of the surplus. Another person can look modest on an income statement while quietly building substantial wealth through assets that keep compounding.

\section{Why people accept money}

People accept money because it sits inside a web of trust and enforcement:

\begin{itemize}
  \item taxes are due in it,
  \item contracts are written in it,
  \item banks and payment rails settle in it,
  \item and everyone expects everyone else to accept it.
\end{itemize}

Money, then, is not only an object. It is also a legal and social arrangement.

\chapter{Where money comes from}

This question contains two different levels of analysis: the macro system and the household.

\section{System-level creation}

The Bank of England's explainer on money creation remains one of the clearest official public summaries of the modern process: most money in the economy is created by commercial banks when they make loans. That alone clears up a common misunderstanding. Banks do not simply lend out pre-existing piles of household deposits in a mechanical way. Lending creates deposits.

\begin{notebox}
\textbf{Operational distinction}\par
\textbf{Money creation} is about the architecture of the monetary system.\par
\textbf{Money earning} is about who captures value inside that system.
\end{notebox}

At the macro level, the main channels are:

\begin{center}
\begin{tabularx}{\textwidth}{>{\raggedright\arraybackslash}p{0.25\textwidth}X>{\raggedright\arraybackslash}p{0.22\textwidth}}
\toprule
\textbf{Channel} & \textbf{Mechanism} & \textbf{Practical implication} \\
\midrule
Commercial bank lending & New loans create new deposits. & Broad money expands when credit expands. \\
Central bank liabilities & The central bank issues reserves and currency. & This anchors settlement, liquidity, and monetary policy. \\
Government spending and transfers & Public spending injects purchasing power into firms and households. & Public budgets shape demand, safety nets, and infrastructure. \\
External inflows & Exports, remittances, and capital inflows add purchasing power. & Open economies can receive strong support from outside demand and capital. \\
\bottomrule
\end{tabularx}
\end{center}

\section{Household-level acquisition}

At the personal level, money usually arrives through:

\begin{itemize}
  \item wages and salaries,
  \item business revenue,
  \item freelance or service income,
  \item investment income,
  \item asset sales,
  \item public transfers,
  \item gifts and inheritances,
  \item and sometimes debt.
\end{itemize}

That list is useful because it forces a practical question: which channels are one-time, and which can become repeatable without requiring the same effort every time?

\section{Channel map: from system credit to household outcomes}

A common error is to track only one layer (for example, policy rates) and ignore access, flow, and stress transmission.

\begin{center}
\small
\begin{tabularx}{\textwidth}{>{\raggedright\arraybackslash}p{0.18\textwidth}>{\raggedright\arraybackslash}p{0.22\textwidth}>{\raggedright\arraybackslash}p{0.24\textwidth}X}
\toprule
\textbf{Layer} & \textbf{What to observe} & \textbf{Primary source} & \textbf{Why it matters} \\
\midrule
Money and balance-sheet expansion & Broad money and sector balance-sheet growth. & Fed H.6 + Fed Z.1 & Indicates whether system liquidity is expanding or contracting. \\
Credit access conditions & Lending standards and approval friction. & Fed SLOOS & Tightening can block ownership entry before headline stress appears. \\
Mortgage entry and borrower mix & Loan-level mortgage distribution and origination composition. & HMDA modified LAR + CFPB HMDA release notes & Helps separate incumbent gains from new-owner entry. \\
Household stress realization & Delinquency transitions and debt-burden outcomes. & NY Fed Household Debt and Credit + Fed DSR/FOR & Confirms whether tighter access is translating into balance-sheet damage. \\
Mortgage performance detail & Origination, payment, and performance structure. & FHFA NMDB & Adds granular context to leverage quality and fragility. \\
Small-business credit conditions & Approval, denial, and financing constraints for firms under 500 employees. & Fed SBCS + CFPB Small Business Lending & Tracks business-equity ownership channels beyond wage income. \\
\bottomrule
\end{tabularx}
\normalsize
\end{center}

\begin{notebox}
\textbf{Practical read}\par
If this full sequence is improving, ownership entry conditions are usually broadening. If access tightens while burden and delinquencies rise, prioritize liquidity and resilience before adding leverage.
\end{notebox}

\chapter{What wealth is}

The accounting identity is simple:

\begin{center}
\fbox{\parbox{0.72\textwidth}{\centering \large \textbf{Wealth = Assets - Liabilities}}}
\end{center}

But life is more interesting than a single equation. Wealth is also the ability of your assets, skills, and systems to keep generating future options.

\section{The six capitals view}

\begin{center}
\begin{tabularx}{\textwidth}{>{\raggedright\arraybackslash}p{0.23\textwidth}X X}
\toprule
\textbf{Capital type} & \textbf{Examples} & \textbf{Wealth effect} \\
\midrule
Financial capital & Cash, stocks, bonds, funds & Generates optionality, interest, dividends, liquidity. \\
Human capital & Skills, health, judgment, reputation & Raises earning power and resilience. \\
Business capital & Equity, customer relationships, retained earnings & Scales profits beyond a single worker's hours. \\
Intellectual capital & Software, research, brand, process know-how & Makes high-margin repeatable output possible. \\
Social capital & Trust, network quality, distribution access & Improves deal flow and lowers friction. \\
Physical capital & Property, equipment, infrastructure & Supports rent, production, and durability. \\
\bottomrule
\end{tabularx}
\end{center}

The World Bank's \textit{Changing Wealth of Nations 2024} report is helpful here because it insists that wealth cannot be reduced to GDP. A country may grow output while depleting the underlying asset base that future growth depends on. The same is true for a household or a business. Revenue alone can hide depletion.

\section{Income versus ownership}

Labor income is necessary. Ownership income changes the slope of the curve.

Workers are paid once for a unit of labor. Owners may get paid repeatedly from the same asset:

\begin{itemize}
  \item dividends from equity,
  \item rent from property,
  \item royalties from intellectual property,
  \item subscription revenue from software,
  \item licensing income from systems and templates.
\end{itemize}

This is one reason wealth inequality persists. The more income is tied to assets rather than hours, the easier it becomes to compound.

\chapter{Why wealth is uneven}

Any honest discussion of wealth has to hold two truths at once:

\begin{itemize}
  \item personal decisions matter,
  \item and starting conditions matter.
\end{itemize}

\section{Structural differences}

Large wealth gaps are shaped by:

\begin{itemize}
  \item family wealth and inheritance,
  \item education quality,
  \item health,
  \item geography,
  \item institutional stability,
  \item discrimination,
  \item access to credit,
  \item and exposure to inflation, violence, or policy failure.
\end{itemize}

These are not excuses. They are real constraints that change what is feasible, how long the path takes, and how much risk a person can absorb.

\section{Strategic differences}

Even among people with similar beginnings, strong divergences emerge from:

\begin{itemize}
  \item savings rate,
  \item skill scarcity,
  \item ownership share,
  \item concentration versus diversification,
  \item time horizon,
  \item decision quality under uncertainty,
  \item and the ability to avoid ruin.
\end{itemize}

\section{Luck and path dependence}

Luck matters more than people like to admit:

\begin{itemize}
  \item entering a rising industry,
  \item finding an extraordinary mentor,
  \item meeting a cofounder,
  \item starting in the right city,
  \item surviving a crisis that eliminates competitors.
\end{itemize}

The right conclusion is not fatalism. It is to build a process that benefits from good luck when it appears and survives bad luck when it arrives.

\section{Access-and-resilience scoreboard (quarterly)}

To keep this chapter operational, track one compact scoreboard rather than relying on narrative instinct.

\begin{center}
\small
\begin{tabularx}{\textwidth}{>{\raggedright\arraybackslash}p{0.18\textwidth}>{\raggedright\arraybackslash}p{0.20\textwidth}>{\raggedright\arraybackslash}p{0.24\textwidth}X}
\toprule
\textbf{Layer} & \textbf{Signal to track} & \textbf{Primary source} & \textbf{Decision use} \\
\midrule
System stress & Broad market stress pulse. & OFR Financial Stress Index (\url{financialresearch.gov/financial-stress-index}) & Falling stress with stable access supports gradual risk re-entry. \\
Bank resilience & Vulnerability narrative and channel-level stress commentary. & Fed Financial Stability Report (\url{federalreserve.gov/publications/financial-stability-report.htm}) & Improving resilience does not automatically mean broad credit access. \\
Credit standards & Lender tightening/easing by segment. & Fed SLOOS (\url{federalreserve.gov/data/sloos.htm}) & Tightening is an early warning for entrant access constraints. \\
Entrant credit access & Small-business and community-lending distribution. & FFIEC CRA Data Products (\url{ffiec.gov/data/cra/data-products}) & Weak entrant credit despite lower stress implies unequal transmission. \\
Ownership-entry pipeline & Firm entry/exit and cohort dynamics. & Census BDS API (\url{census.gov/topics/business-economy/dynamics/data/api.html}) & Strong formation with weak persistence signals churn, not durable ownership gains. \\
Household stress realization & Delinquency transitions and debt-burden outcomes. & NY Fed Household Debt + Fed DSR/FOR & Confirms whether access and burden are becoming balance-sheet damage. \\
\bottomrule
\end{tabularx}
\normalsize
\end{center}

\begin{notebox}
\textbf{Asymmetric cycle warning}\par
If system stress improves but entrant access and household outcomes do not, treat the cycle as asymmetric: resilience for incumbents, friction for new owners.
\end{notebox}

\chapter{How wealth is built in practice}

Wealth typically grows through a stack of six engines plus one execution loop rather than a single trick.

\section{Engine 1: increase earning power}

For most people, the highest-return early move is increasing the value of their work. Rare skills, reliable execution, communication ability, sales ability, and domain judgment often beat portfolio optimization in the beginning.

\section{Engine 2: create margin}

If income rises while spending rises just as fast, wealth stalls. Margin is what lets a future exist. That usually means:

\begin{itemize}
  \item automating transfers,
  \item keeping fixed costs sane,
  \item avoiding chronic high-interest debt,
  \item and making lifestyle inflation a choice rather than an instinct.
\end{itemize}

\section{Engine 3: buy productive assets}

Durable wealth tends to come from ownership of productive assets such as diversified equity, profitable businesses, intellectual property, or cash-flowing property bought at disciplined prices.

\section{Engine 4: use leverage carefully}

Not all leverage is debt. In modern work, the most useful leverage is often:

\begin{itemize}
  \item code,
  \item media,
  \item audience,
  \item systems,
  \item and reusable products.
\end{itemize}

Financial leverage can help, but it punishes mistakes brutally. Informational, technical, and distribution leverage are often safer for builders than borrowed money.

\section{Engine 5: stay in the game}

Compounding only works if you avoid ruin. Insurance, diversification, liquidity buffers, and modest position sizing are not boring side notes. They are the conditions that allow compounding to survive.

\begin{signalbox}
\textbf{A practical formula}\par
Wealth growth = (earning power x savings rate x years x compounding x ownership share) - ruin
\end{signalbox}

\section{Ownership entry execution loop (quarterly)}

Use this loop when deciding whether to accelerate ownership (housing, business equity, or market risk) or stay defensive.

\begin{enumerate}
  \item \textbf{Access check}: are standards tightening, and are entrant channels improving or worsening? (Fed SLOOS, HMDA, SBCS)
  \item \textbf{Burden check}: is payment burden rising faster than real income? (Fed DSR/FOR, BEA Personal Income, BLS CPI)
  \item \textbf{Stress check}: are expectations and delinquency transitions worsening together? (NY Fed SCE, NY Fed Household Debt and Credit)
  \item \textbf{Action check}: if two or more checks are red, prioritize liquidity and lower fixed commitments before adding leverage.
\end{enumerate}

\begin{center}
\small
\begin{tabularx}{\textwidth}{>{\raggedright\arraybackslash}p{0.24\textwidth}>{\raggedright\arraybackslash}p{0.34\textwidth}X}
\toprule
\textbf{Decision state} & \textbf{Signal pattern} & \textbf{Practical move} \\
\midrule
Defensive & Tightening access, rising burden, worsening transitions. & Delay large leverage, increase runway, and protect downside. \\
Neutral & Mixed signals with no broad stress acceleration. & Continue gradual ownership accumulation and avoid concentrated bets. \\
Offensive (controlled) & Stable access, manageable burden, contained transitions. & Add ownership exposure in staged tranches with fixed risk limits. \\
\bottomrule
\end{tabularx}
\normalsize
\end{center}

\section{Scenario-ready leverage stress test (10-minute version)}

Before increasing leverage, run one scenario stress check using supervisory-style shocks:

\begin{itemize}
  \item Scenario anchor: Fed 2026 stress paths (\url{federalreserve.gov/publications/2026-stress-test-scenarios.htm}).
  \item Burden anchor: Fed DSR/FOR (\url{federalreserve.gov/releases/dsr/}).
  \item Household stress anchor: NY Fed Household Debt and Credit (\url{newyorkfed.org/microeconomics/hhdc/background}).
  \item Asset-side anchor: FHFA HPI (\url{fhfa.gov/house-price-index}) plus BLS CPI (\url{bls.gov/cpi/home.htm}).
\end{itemize}

\begin{center}
\small
\begin{tabularx}{\textwidth}{>{\raggedright\arraybackslash}p{0.19\textwidth}>{\raggedright\arraybackslash}p{0.24\textwidth}>{\raggedright\arraybackslash}p{0.24\textwidth}X}
\toprule
\textbf{Test block} & \textbf{Pass condition} & \textbf{Fail condition} & \textbf{Action if fail} \\
\midrule
Income shock tolerance & Cash-flow plan survives unemployment/income-hit assumptions. & Negative monthly cash flow for multiple quarters. & Reduce fixed commitments and delay new leverage. \\
Payment burden tolerance & Debt-service share remains within pre-set safe range. & Burden rises above your red-line threshold. & Deleverage first, then reassess. \\
Asset drawdown tolerance & No forced sale under house/equity drawdown. & Margin call, forced refi, or liquidity break. & Increase liquidity buffer and reduce concentration. \\
Inflation-adjusted resilience & Real income holds up versus fixed obligations. & Real income erosion plus rising debt burden. & Prioritize pricing power, skill/wage upgrades, and lower fixed costs. \\
\bottomrule
\end{tabularx}
\normalsize
\end{center}

\begin{signalbox}
\textbf{Rule}\par
Do not scale leverage only because stress indicators are quiet; scale only after your own balance sheet passes this stress check.
\end{signalbox}

\section{Distinction checklist before major money decisions}

Use this checklist before changing jobs, taking debt, buying assets, or changing long-term allocation.

\begin{center}
\small
\begin{tabularx}{\textwidth}{>{\raggedright\arraybackslash}p{0.24\textwidth}>{\raggedright\arraybackslash}p{0.28\textwidth}X}
\toprule
\textbf{Distinction} & \textbf{Common mistake} & \textbf{Better prompt} \\
\midrule
Income flow vs wealth stock & Treating a raise as wealth creation. & How much of this new income turns into owned assets after 12 months? \\
Nominal gain vs real gain & Ignoring inflation and tax drag. & What is the after-tax, after-inflation result? \\
Average outcome vs median outcome & Copying advice built for top deciles. & What happens at the median household, not only the mean? \\
Volatility vs ruin & Accepting hidden blow-up risk. & Can this plan survive a bad 2-3 year sequence? \\
Productive leverage vs fragile leverage & Using debt for status consumption. & Does this leverage increase durable cash flow or only fixed obligations? \\
Money creation vs money distribution & Assuming liquidity reaches everyone equally. & Who gets the new credit first, and who carries the lag? \\
Asset inflation vs goods inflation & Confusing portfolio gains with purchasing power. & Did real spending power rise, or only mark-to-market value? \\
\bottomrule
\end{tabularx}
\normalsize
\end{center}

\section{Evidence ladder for financial claims}

When a claim sounds attractive, rank it before acting:

\begin{enumerate}
  \item \textbf{Anecdote}: one person, one cycle, low reliability.
  \item \textbf{Backtest or narrative}: useful for hypotheses, not final decisions.
  \item \textbf{Cross-sectional evidence}: multiple households or countries, still fragile to omitted variables.
  \item \textbf{Long-run panel evidence}: better for stress-testing across regimes.
  \item \textbf{Official statistical releases}: baseline for macro and household reality checks.
\end{enumerate}

Practical rule: only take concentrated risk when claim quality is at least level 4 and downside remains survivable.

\section{Household stress dashboard (monthly or quarterly)}

If the goal is durable wealth, track stress signals before visible balance-sheet damage appears.

\begin{center}
\small
\begin{tabularx}{\textwidth}{>{\raggedright\arraybackslash}p{0.20\textwidth}>{\raggedright\arraybackslash}p{0.22\textwidth}>{\raggedright\arraybackslash}p{0.20\textwidth}X}
\toprule
\textbf{Signal} & \textbf{What it tracks} & \textbf{Where to check} & \textbf{Practical read} \\
\midrule
Debt-service burden & Fixed payment pressure relative to disposable income. & Federal Reserve DSR/FOR. & Persistent rise means less room for compounding and higher fragility. \\
Expected missed-payment risk & Household expectation of future delinquency. & New York Fed SCE. & Rising expectations often appear before realized defaults. \\
Realized delinquency transitions & Movement from current to 30/60/90+ day delinquency. & New York Fed Household Debt and Credit. & Confirms whether stress is broadening or contained. \\
Difficulty paying bills & Self-reported financial strain and liquidity stress. & CFPB Making Ends Meet. & Captures household pressure not always visible in credit files. \\
Real income momentum & Income growth net of inflation. & BEA Personal Income + BLS CPI. & Weak real income with high debt service is a high-risk mix. \\
\bottomrule
\end{tabularx}
\normalsize
\end{center}

Use order of operations:

\begin{enumerate}
  \item Track expectations (SCE) and debt-service burden (DSR/FOR) as early signals.
  \item Confirm with realized delinquency transitions (NY Fed Household Debt).
  \item Adjust leverage, liquidity buffers, and fixed-cost commitments before stress compounds.
\end{enumerate}

\section{Credit-conditions transmission check (monthly or quarterly)}

A recurring mistake is to watch only interest rates while ignoring access conditions.

Use this transmission sequence:

\begin{enumerate}
  \item \textbf{Access}: are banks tightening standards? (Fed SLOOS)
  \item \textbf{Flow}: are household credit originations and inquiries slowing? (CFPB Consumer Credit Trends)
  \item \textbf{Outcome}: are debt burden and delinquencies worsening? (BIS DSR + NY Fed Household Debt)
\end{enumerate}

\begin{center}
\small
\begin{tabularx}{\textwidth}{>{\raggedright\arraybackslash}p{0.18\textwidth}>{\raggedright\arraybackslash}p{0.20\textwidth}>{\raggedright\arraybackslash}p{0.22\textwidth}X}
\toprule
\textbf{Signal layer} & \textbf{What it tracks} & \textbf{Where to check} & \textbf{Practical read} \\
\midrule
Credit standards & Lender risk appetite and approval friction. & Federal Reserve SLOOS. & Tighter standards can reduce access before policy-rate effects appear in balances. \\
Credit flow & Origination and inquiry momentum by product and risk slice. & CFPB Consumer Credit Trends. & Flow slowdown after tightening suggests transmission is active. \\
Debt burden & Debt-service pressure relative to income. & BIS DSR + Fed DSR/FOR. & Persistent burden rise narrows margin for error and weakens compounding capacity. \\
Realized stress & Transitions into delinquency buckets. & NY Fed Household Debt and Credit. & Confirms whether pressure is becoming observable balance-sheet damage. \\
Household resilience & Bill stress, shock capacity, and coping behavior. & Fed SHED. & Cross-checks whether macro credit signals match lived household conditions. \\
\bottomrule
\end{tabularx}
\normalsize
\end{center}

If standards tighten but flow does not slow, transmission may be delayed or composition-driven. If flow slows and delinquencies rise together, prioritize liquidity and deleveraging over return-chasing.

\chapter{Who can get wealth and what should they do first}

Almost anyone can improve their financial condition. Not everyone should use the same strategy. A useful sequence is:

\begin{center}
\begin{tabularx}{\textwidth}{>{\raggedright\arraybackslash}p{0.17\textwidth}X X}
\toprule
\textbf{Stage} & \textbf{Goal} & \textbf{Best first move} \\
\midrule
Survival & Stop chaos and stabilize cash flow & Eliminate acute leaks, secure income, stop high-cost debt growth. \\
Stability & Build slack & Emergency fund, bill automation, simple budget, insurance basics. \\
Surplus & Create investable cash & Raise income, cut unhelpful recurring costs, automate savings. \\
Ownership & Acquire productive assets & Buy diversified assets or build small owned products. \\
Autonomy & Make work more optional & Increase recurring cash flow and improve system resilience. \\
\bottomrule
\end{tabularx}
\end{center}

The first goal is not status. It is margin. The second is ownership.

\chapter{Can a thing get money?}

Yes. A company, property, fund, or software product can receive cash flows because legal and economic systems assign rights to those structures.

\begin{itemize}
  \item A rental property receives rent.
  \item A company receives revenue.
  \item A software product receives subscriptions.
  \item A fund receives dividends and interest from underlying holdings.
\end{itemize}

But assets do not enjoy the result. Owners, creditors, and beneficiaries ultimately capture the economic benefit. This is why the deepest practical wealth question is not only ``How do I make money?'' but ``What do I own that keeps producing after I stop pushing?''

\chapter{A 90-day operating system}

\section{Days 1-30: map reality}

\begin{itemize}
  \item Calculate net worth.
  \item List every debt and interest rate.
  \item Track spending for one month.
  \item Build a one-page personal balance sheet.
  \item Identify one skill with the highest probability of lifting income.
\end{itemize}

\section{Days 31-60: create margin}

\begin{itemize}
  \item Cut one recurring cost that adds little value.
  \item Automate a transfer to savings or investment.
  \item Reduce the highest-friction debt.
  \item Build a starter emergency fund.
  \item Launch one simple side-income offer.
\end{itemize}

\section{Days 61-90: buy or build assets}

\begin{itemize}
  \item Fund a diversified investment account.
  \item Build one reusable digital asset: template, guide, calculator, lesson, or workflow.
  \item Set a weekly review ritual.
  \item Choose one long game and commit to it for a year.
\end{itemize}

\chapter{Resource map}

\section{Books}

\begin{longtable}{>{\raggedright\arraybackslash}p{0.28\textwidth}p{0.22\textwidth}p{0.38\textwidth}}
\toprule
\textbf{Book} & \textbf{Best for} & \textbf{Why it matters} \\
\midrule
\endhead
\textit{The Psychology of Money} by Morgan Housel & behavior & A clear guide to how temperament shapes financial outcomes. \\
\textit{The Intelligent Investor} by Benjamin Graham & investing discipline & Classic value-investing framework and defensive mindset. \\
\textit{A Random Walk Down Wall Street} by Burton G. Malkiel & index investing & Strong baseline case for low-cost evidence-aware investing. \\
\textit{I Will Teach You to Be Rich} by Ramit Sethi & cash-flow systems & Practical setup for accounts, automation, spending, and investing. \\
\textit{The Millionaire Next Door} by Thomas J. Stanley and William D. Danko & lifestyle design & Useful reminder that quiet behavior often beats visible consumption. \\
\textit{The Ascent of Money} by Niall Ferguson & financial history & Places modern finance inside a long historical arc. \\
\textit{The Simple Path to Wealth} by J.L. Collins & simple accumulation & Clear, low-friction long-term compounding approach. \\
\textit{Capital in the Twenty-First Century} by Thomas Piketty & inequality & Essential if you want to understand concentration of wealth at scale. \\
\bottomrule
\end{longtable}

\section{Courses and tutorials}

\begin{longtable}{>{\raggedright\arraybackslash}p{0.32\textwidth}p{0.22\textwidth}p{0.32\textwidth}}
\toprule
\textbf{Resource} & \textbf{Format} & \textbf{Use} \\
\midrule
\endhead
Yale / Coursera: \textit{Financial Markets} & online course & Markets, risk, behavioral finance, and institutional basics. \\
Open Yale Courses: \textit{Financial Theory} & free lectures & Finance as part of the wider economic system. \\
MIT OCW: \textit{Blockchain and Money} & free graduate course & Money, payments, crypto, and regulation. \\
MIT OCW: \textit{Principles of Microeconomics} & free course & Incentives, trade-offs, pricing, and firm behavior. \\
OpenStax: \textit{Principles of Economics 3e} & free textbook & Broad economics foundation. \\
Khan Academy: \textit{Personal Finance} & free course & Saving, debt, taxes, insurance, and investing basics. \\
Investor.gov calculators & practical tools & Compound-interest and savings-goal modeling. \\
CFPB: \textit{Your Money, Your Goals} & toolkit & Goals, bills, debt, credit, and consumer decision support. \\
Federal Reserve Education: \textit{Making Personal Finance Decisions} & curriculum & Twenty lessons grounded in economic theory. \\
\bottomrule
\end{longtable}

\section{Online repositories}

\begin{itemize}
  \item \href{https://github.com/joshjluo/personal-finance-guide}{joshjluo/personal-finance-guide}
  \item \href{https://github.com/philschatz/economics-book}{philschatz/economics-book}
  \item \href{https://github.com/VladZn/gb-personal-finance}{VladZn/gb-personal-finance}
  \item \href{https://github.com/brandonhimpfen/awesome-finance}{brandonhimpfen/awesome-finance}
  \item \href{https://github.com/georgezouq/awesome-ai-in-finance}{georgezouq/awesome-ai-in-finance}
\end{itemize}

\section{Official data and references}

\begin{longtable}{>{\raggedright\arraybackslash}p{0.28\textwidth}p{0.34\textwidth}p{0.26\textwidth}}
\toprule
\textbf{Source} & \textbf{Use this for} & \textbf{Cadence / scope} \\
\midrule
\endhead
Federal Reserve H.6 & Broad money and component trends. & Weekly/monthly U.S. money stock release. \\
Federal Reserve Z.1 + DFA + EFA & Sector balance sheets, distributional and geographic detail. & Quarterly U.S. flow-of-funds stack. \\
Federal Reserve SCF & Household assets, liabilities, and ownership concentration. & Triennial microdata benchmark. \\
Federal Reserve SLOOS & Lending standards and loan demand by major credit segment. & Quarterly U.S. bank lending survey. \\
Federal Reserve SHED & Household financial well-being, emergency capacity, and payment stress. & Annual U.S. household survey. \\
Federal Reserve DSR/FOR & Household debt-service and fixed-obligation burden. & Quarterly U.S. household burden release. \\
New York Fed Household Debt and Credit & Delinquency transitions and debt composition by household slices. & Quarterly U.S. credit panel. \\
New York Fed SCE & Inflation, labor, credit access, and expected delinquency signals. & Monthly U.S. household expectations survey. \\
CFPB Consumer Credit Trends & Consumer credit flow by product/risk group and geography. & Monthly U.S. dashboard updates. \\
CFPB Making Ends Meet & Direct household stress and bill-payment strain indicators. & Annual survey reports with data files. \\
BEA Personal Income and Outlays & Income, consumption, and personal saving path. & Monthly U.S. national accounts release. \\
BEA Distribution of Personal Income & Distributional disposable income and inequality decomposition. & Annual distributional national-accounts release. \\
BLS CPI + CEX & Inflation pressure and household spending structure. & Monthly CPI and annual spending detail. \\
FHFA House Price Index & Housing wealth regime and regional price dynamics. & Monthly/quarterly U.S. house-price indices. \\
FHFA National Mortgage Database (NMDB) & Mortgage origination, performance, and outstanding-balance structure. & Release-based U.S. mortgage sample and aggregate statistics. \\
FFIEC HMDA modified LAR + CFPB HMDA release notes & Loan-level mortgage distribution, originations, and borrower mix context. & Annual HMDA filing cycle with current release notes. \\
Federal Reserve Small Business Credit Survey (SBCS) & Credit access, denial, and financing conditions for firms under 500 employees. & Annual survey reports with subgroup cut tables. \\
CFPB Small Business Lending (Section 1071) & Filing and publication surface for small-business lending disclosures. & Rule-driven filing/publishing cadence with specification updates. \\
BIS Debt Service Ratios (DSR) & Cross-country debt burden and early-warning leverage context. & Quarterly international debt-service statistics. \\
U.S. Census wealth tables (SIPP) & Household wealth and debt distribution in public-use tables. & Annual SIPP-based wealth publication. \\
U.S. Census CPS income/inequality tables & Household, family, and person income distribution baselines. & Annual CPS ASEC update cycle. \\
IRS SOI Publication 1304 & AGI and tax-share distribution cross-check from filed returns. & Annual individual return statistics report. \\
WID + WIID + OECD IDD/WDD & Cross-country inequality and distribution comparisons. & Global inequality and OECD harmonized datasets. \\
World Bank CWON + PIP + Findex + IDS & Comprehensive wealth, poverty, inclusion, and debt context. & Global development and debt datasets. \\
IMF WEO + GDD + BIS data portal + ECB CES & Macro regime, debt cycle, and expectations-sensitive cross-country benchmarks. & Global macro-financial and expectations references. \\
SEC EDGAR API docs + FRED/ALFRED & Reproducible filings ingestion and revision-aware macro pulls. & API docs plus vintage-aware archive layer. \\
FRASER historical archive & Primary-source monetary and financial history documents. & Long-run policy and publication archive. \\
\bottomrule
\end{longtable}

\chapter{Source notes}

Primary references checked for this edition on 2026-04-05:

\begin{itemize}
  \item Federal Reserve releases: \href{https://www.federalreserve.gov/releases/h6/current/default.htm}{H.6}, \href{https://www.federalreserve.gov/releases/z1/}{Z.1}, \href{https://www.federalreserve.gov/releases/z1/dataviz/dfa/index.html}{DFA}, \href{https://www.federalreserve.gov/releases/efa/enhanced-financial-accounts.htm}{EFA}, \href{https://www.federalreserve.gov/econres/scfindex.htm}{SCF}, and \href{https://www.federalreserve.gov/releases/dsr/}{DSR/FOR}
  \item Federal Reserve surveys: \href{https://www.federalreserve.gov/data/sloos.htm}{SLOOS} and \href{https://www.federalreserve.gov/publications/report-economic-well-being-us-households.htm}{SHED}
  \item Federal Reserve Bank of New York, \href{https://www.newyorkfed.org/microeconomics/hhdc}{Household Debt and Credit} and \href{https://www.newyorkfed.org/microeconomics/sce}{Survey of Consumer Expectations}
  \item U.S. BEA, \href{https://www.bea.gov/data/income-saving/personal-income}{Personal Income and Outlays} and \href{https://www.bea.gov/data/special-topics/distribution-of-personal-income}{Distribution of Personal Income}
  \item U.S. BLS, \href{https://www.bls.gov/cpi/home.htm}{CPI} and \href{https://www.bls.gov/cex/}{Consumer Expenditure Surveys}
  \item U.S. Census, \href{https://www.census.gov/topics/income-poverty/wealth/data/tables.html}{Wealth and Asset Ownership tables (SIPP)}, \href{https://www.census.gov/topics/income-poverty/income-inequality/data/data-tables/cps-data-tables.html}{CPS income/inequality tables}, and \href{https://www2.census.gov/library/publications/2025/demo/p70br-211.pdf}{P70BR-211}
  \item U.S. FHFA, \href{https://www.fhfa.gov/house-price-index}{House Price Index datasets} and \href{https://www.fhfa.gov/data/nmdb}{National Mortgage Database (NMDB)}
  \item IRS SOI, \href{https://www.irs.gov/statistics/soi-tax-stats-individual-income-tax-returns-complete-report-publication-1304}{Publication 1304} and \href{https://www.irs.gov/statistics/soi-tax-stats-irs-data-book}{IRS Data Book}
  \item Bank of England, \href{https://www.bankofengland.co.uk/quarterly-bulletin/2014/q1/money-creation-in-the-modern-economy}{Money creation in the modern economy}
  \item IMF, \href{https://data.imf.org/Datasets/WEO}{World Economic Outlook} and \href{https://data.imf.org/en/datasets/IMF.FAD%3AGDD}{Global Debt Database}
  \item BIS, \href{https://data.bis.org/}{Data Portal} and \href{https://data.bis.org/topics/DSR}{Debt Service Ratios}
  \item ECB, \href{https://www.ecb.europa.eu/stats/ecb_surveys/consumer_exp_survey/html/index.en.html}{Consumer Expectations Survey}
  \item World Bank, \href{https://www.worldbank.org/en/publication/the-changing-wealth-of-nations}{The Changing Wealth of Nations 2024}, \href{https://www.worldbank.org/en/publication/globalfindex/download-data}{Global Findex}, \href{https://datacatalog.worldbank.org/search/dataset/0038020/poverty-and-equity-database-poverty-and-inequality-platform}{PIP}, and \href{https://www.worldbank.org/en/programs/debt-statistics/ids}{IDS}
  \item \href{https://wid.world/}{WID.world}, \href{https://www.wider.unu.edu/database/world-income-inequality-database-wiid}{WIID}, and \href{https://www.oecd.org/en/data/datasets/income-and-wealth-distribution-database.html}{OECD IDD/WDD}
  \item HMDA sources, \href{https://ffiec.cfpb.gov/data-publication/modified-lar}{FFIEC modified LAR} and \href{https://www.consumerfinance.gov/about-us/newsroom/2025-hmda-data-on-mortgage-lending-now-available/}{CFPB 2025 HMDA release notes}
  \item Federal Reserve Small Business Credit Survey, \href{https://www.fedsmallbusiness.org/reports/survey/}{all survey years reports}
  \item CFPB, \href{https://www.consumerfinance.gov/data-research/consumer-credit-trends/}{Consumer Credit Trends}, \href{https://www.consumerfinance.gov/data-research/making-ends-meet-survey-data/}{Making Ends Meet survey data}, \href{https://www.consumerfinance.gov/data-research/small-business-lending/}{Small Business Lending}, and \href{https://www.consumerfinance.gov/your-money-your-goals/}{Your Money, Your Goals}
  \item SEC, \href{https://www.sec.gov/edgar/sec-api-documentation}{EDGAR API documentation}
  \item St. Louis Fed, \href{https://fred.stlouisfed.org/docs/api/fred/fred/}{FRED API docs}, \href{https://alfred.stlouisfed.org/help/downloaddata}{ALFRED}, and \href{https://fraser.stlouisfed.org/}{FRASER}
  \item Yale / Coursera, \href{https://www.coursera.org/learn/financial-markets-global}{Financial Markets}
  \item Open Yale Courses, \href{https://oyc.yale.edu/economics/econ-251}{Financial Theory}
  \item MIT OpenCourseWare, \href{https://ocw.mit.edu/courses/15-s12-blockchain-and-money-fall-2018/}{Blockchain and Money}
  \item OpenStax, \href{https://openstax.org/books/principles-economics-3e/pages/preface}{Principles of Economics 3e}
  \item Khan Academy, \href{https://www.khanacademy.org/college-careers-more/personal-finance}{Personal Finance}
  \item Investor.gov, \href{https://www.investor.gov/use-financial-tools-and-calculators}{Use Financial Tools and Calculators}
  \item Federal Reserve Education, \href{https://www.federalreserveeducation.org/teaching-resources/personal-finance/curriculum/making-personal-finance-decisions-curriculum}{Making Personal Finance Decisions Curriculum}
\end{itemize}

\backmatter

\begin{signalbox}
\textbf{Final takeaway}\par
The most useful wealth strategy is usually not flashy: earn better, spend with intention, own productive things, compound patiently, and avoid decisions that remove you from the game.
\end{signalbox}

\end{document}
